\chapter{Normativa y ruta de autorizaciones}
\section{Aplicación de CNE Suministro y Utilización}
El CNE Suministro 2011 regula los aspectos de suministro, estaciones, líneas, conductores, distancias, maniobra y seguridad de la infraestructura. El CNE Utilización, aprobado por RM 037-2006 con la modificación RM 175-2008, se aplica a las instalaciones de utilización que corresponden, incluidos los circuitos auxiliares y receptores. No se emplea una regla residencial para sustituir el diseño de una estación generadora ni se atribuye al CNE una potencia por vivienda que procede de la norma rural de 2023.

La matriz del anexo A registra disposición, criterio, evidencia y estado abierto. Se revisaron particularmente CNE Suministro 036 (puesta a tierra), 125 (espacios), 130 (sobrevelocidad), 132 (neutros) y reglas de líneas 230, 232 y 234. CNE Utilización 030-004, 050-104, 050-106, sección 060 y tablas de ampacidad se consultan junto a instalación real y marcado de equipos. La existencia de una referencia bibliográfica no acredita cumplimiento de todas sus reglas.

\section{Cambios normativos importantes para esta revisión}
La RD 108-2023-MINEM/DGE sustituye las antiguas bases rurales de redes secundarias aprobadas por RD 031-2003. Sus valores de calificación y simultaneidad modifican sustancialmente la demanda, por lo que se recalculan ambos y se conserva el ámbito académico de aplicación. La E.030 fue modificada en 2026; no se reutiliza una zonificación histórica sin verificar el anexo distrital actual. Se identifica además la disposición transitoria RM 217-2026, cuya aplicación a una obra real depende de las condiciones de tramitación.

La ficha oficial de IEC 60909-0 identifica una edición 2026 que reemplaza la de 2016. IEC 60034-1:2026 y IEC 60076-1:2011 sirven como referencias de especificación de máquinas y transformadores. La consulta disponible de estas normas internacionales es de catálogo, no de sus textos completos pagados. No se inventan subcláusulas ni se certifican diseños conforme a ellas. Antes de una contratación se deben adquirir las ediciones pertinentes y precisar ensayos, pérdidas, elevación de temperatura, cortocircuito, aislamiento y funcionamiento a la altitud del sitio.

\section{Agua, ambiente, estructuras y seguridad}
Ley 29338 y su Reglamento distinguen disponibilidad, derecho de uso, intervención en cauce y autorización de obras. Se revisan, entre otros, artículos 44, 47 y 53 de la ley para uso y licencia; 74 para faja marginal; 81 para evaluación ambiental; y 104 para obras. La modificación DS 023-2014-MINAGRI se consulta junto a los artículos reglamentarios pertinentes. El procedimiento ANA de una ejecución real deberá verificarse con el TUPA y las disposiciones vigentes; no se declara lectura integral de anexos no recuperados de RJ 007-2015.

El caudal ecológico se desarrolla según RJ 267-2019, numeral 6: metodología dependiente de categoría y coordinación con la autoridad. El porcentaje de estudio utilizado no es una obligación legal ni un caudal aprobado. Deben examinarse usos poblacionales, agrarios, energéticos y restricciones del tramo, así como disponibilidad en estiaje y retorno del agua.

El DS 014-2019-EM, junto a su fe de erratas, regula protección ambiental de actividades eléctricas. La categoría no se determina solamente porque la potencia sea pequeña. El tipo de aprovechamiento, embalse, longitud de línea, ubicación y sensibilidad ambiental pueden modificar el instrumento exigible. Se revisan criterios del anexo correspondiente sin afirmar una DIA automática. El DS 011-2026-EM se refiere a adecuación de actividades en curso; no constituye certificación ambiental de una nueva ampliación. Debe distinguirse situación legal de la instalación existente, rehabilitación y obras nuevas.

Para obras civiles se aplicará el RNE pertinente: E.030 para sismo, E.050 para suelos y cimentaciones, E.060 para concreto armado, además de reglas de arquitectura, accesos y seguridad que correspondan. El diseño estructural definitivo requiere geotecnia y modelos de carga. Ley 29783 y el reglamento de seguridad eléctrica RM 111-2013 sustentan gestión de riesgos, competencias, consignación, trabajo en altura y planes de emergencia. No se presenta esta memoria estudiantil como documentación firmada por profesionales habilitados para ejecutar.

\section{Obligaciones y recomendaciones técnicas}
Las leyes, reglamentos y códigos nacionales aplicables son obligaciones dentro de su ámbito. ASCE MOP 79, USBR EM03, fichas de fabricante y métodos hidráulicos son apoyos técnicos; solo una especificación contractual o incorporación normativa puede atribuirles otra exigibilidad. Se separan parámetros adoptados de valores medidos y se identifican las ediciones. Donde falta acceso al texto de una norma complementaria, se deja abierta la comprobación correspondiente.

\section{Ruta orientativa para una ejecución real}
Todos los procedimientos están pendientes; no hay permisos obtenidos.
\begin{enumerate}
\item Identificar titular, estado legal y bienes de la central existente. Levantar instalaciones y verificar propiedad, acceso, servidumbres y ámbito de prestación. La mención Municipalidad/Junta en antecedentes no acredita titularidad.
\item Consultar con ANA/AAA competente disponibilidad, derechos concurrentes, caudal ecológico, faja marginal y autorizaciones necesarias. Tramitar los procedimientos aplicables y la licencia de uso; no atribuir la licencia de Phara a Limbani.
\item Definir con la autoridad ambiental el instrumento y la categoría para rehabilitación, ampliación y distribución. Incorporar línea, combustible, sedimentos, ruido, afectación de cauce y patrimonio. Evaluar permisos forestales, arqueológicos y prediales según afectaciones reales.
\item Verificar competencia MINEM/Gobierno Regional de Puno y el régimen eléctrico vigente. La modificación introducida por DL 1002 considera concesión definitiva para generación hidráulica por encima de 500 kW. El caso estudiado de 600 kW supera ese umbral. También debe evaluarse la distribución de servicio público con demanda superior a 500 kW y la autorización de generación térmica de respaldo superior a 500 kW conforme a los artículos aplicables de LCE. No se conserva una pretendida exención general basada en los 150 kW del borrador.
\item Desarrollar estudios de ingeniería firmados, expediente de obras, autorizaciones de intervención, permisos municipales y de seguridad que procedan. Verificar la condición de sistema eléctrico rural y la normativa específica con la autoridad responsable.
\item Obtener las autorizaciones, ejecutar con control de calidad y realizar pruebas de presión, regulación, protecciones, aislamiento y tierra antes de puesta en servicio. Registrar operación, mantenimiento y responsables.
\end{enumerate}
